\documentclass[12pt, a4paper]{article}
\usepackage[utf8]{inputenc}
\usepackage[spanish]{babel}
\usepackage[T1]{fontenc}
\usepackage{geometry}
\usepackage{amsmath}
\usepackage{listings}
\usepackage{xcolor}
\usepackage{booktabs}
\usepackage{hyperref}
\usepackage{float}
\usepackage{enumitem}
\usepackage{tikz}
\usetikzlibrary{shapes.geometric, arrows.meta, positioning}

\geometry{margin=2.5cm}

\hypersetup{
    colorlinks=true,
    linkcolor=blue!60!black,
    urlcolor=blue!60!black,
    citecolor=blue!60!black,
}

\definecolor{codegreen}{rgb}{0,0.6,0}
\definecolor{codegray}{rgb}{0.5,0.5,0.5}
\definecolor{codepurple}{rgb}{0.58,0,0.82}
\definecolor{backcolour}{rgb}{0.96,0.96,0.96}

\lstdefinestyle{mystyle}{
    backgroundcolor=\color{backcolour},
    commentstyle=\color{codegreen},
    keywordstyle=\color{blue},
    numberstyle=\tiny\color{codegray},
    stringstyle=\color{codepurple},
    basicstyle=\ttfamily\small,
    breaklines=true,
    captionpos=b,
    numbers=left,
    numbersep=5pt,
    showstringspaces=false,
    tabsize=2,
    frame=single,
    framerule=0.4pt,
    rulecolor=\color{black!20},
}
\lstset{style=mystyle}

\begin{document}

% ─── PORTADA ──────────────────────────────
\begin{titlepage}
    \centering
    \vspace*{1cm}
    {\Large \textbf{Universidad del Valle de Guatemala}\\[0.3cm]}
    {\large Facultad de Ingenieria\\[0.1cm]}
    {\large Departamento de Ciencias de la Computacion\\[1.8cm]}

    \rule{\textwidth}{1pt}\\[0.5cm]
    {\LARGE \textbf{Laboratorio 3}\\[0.3cm]}
    {\Large \textbf{Protocolos de Enrutamiento}}\\[0.3cm]
    \rule{\textwidth}{1pt}\\[1.8cm]

    {\large \textbf{CC3067 --- Redes}\\[0.4cm]}
    {\large Semestre II --- 2026\\[2cm]}

    \begin{tabular}{ll}
        \toprule
        \textbf{Integrante}   & \textbf{Carne} \\
        \midrule
        Javier Espana         & 23361          \\
        Roberto Barreda       & 23354          \\
        \bottomrule
    \end{tabular}\\[2cm]

    {\large Guatemala, agosto de 2026}
\end{titlepage}

\tableofcontents
\newpage

% ──────────────────────────────────────────
\section{Descripcion de la Practica}

El Laboratorio 3 de Redes consiste en implementar el protocolo de enrutamiento
\textbf{Link-State} sobre una red simulada mediante sockets TCP. Cada nodo de
la red es un proceso independiente que descubre sus vecinos, intercambia
informacion de estado del enlace y calcula las rutas optimas mediante el
algoritmo de Dijkstra.

La red de prueba se compone de tres parejas de estudiantes (seis integrantes),
cada una responsable de uno o varios nodos. La conectividad distribuida se
logra mediante \textbf{Tailscale}, una VPN P2P que asigna una IP
\texttt{100.x.x.x} fija a cada maquina. Las parejas deben acordar un
\textbf{protocolo comun} para garantizar la interoperabilidad de sus
implementaciones independientes.

La practica se divide en dos partes:
\begin{enumerate}
    \item \textbf{Plano de control}: implementacion del protocolo Link-State
          (HELLO, LSA, Dijkstra, tabla de enrutamiento CSV).
    \item \textbf{Plano de datos}: reenvio salto a salto de mensajes con
          deteccion y correccion de errores Hamming (7,4), mas la interconexion
          y prueba entre todos los grupos.
\end{enumerate}

La topologia acordada para la topologia de seis nodos es la siguiente:

\begin{figure}[H]
\centering
\begin{tikzpicture}[
    router/.style={draw, circle, minimum size=1.1cm, fill=blue!15,
                   font=\small\bfseries},
    host/.style={draw, rectangle, rounded corners=3pt, minimum width=1.5cm,
                 minimum height=0.75cm, fill=green!15, font=\small\bfseries},
    lbl/.style={font=\footnotesize, midway, above, sloped},
]
    \node[router] (R1) at (0,0)     {R1};
    \node[router] (R2) at (3,1.5)   {R2};
    \node[router] (R3) at (0,-2.5)  {R3};
    \node[router] (R4) at (3,-2.5)  {R4};
    \node[router] (R5) at (6,0)     {R5};
    \node[router] (R6) at (6,-2.5)  {R6};

    \node[host] (ATM1)  at (-2.8, 0)   {ATM1};
    \node[host] (BANK1) at (3, -4.5)   {BANK1};

    \draw[thick] (R1) -- node[lbl]{1} (R2);
    \draw[thick] (R2) -- node[lbl]{2} (R3);
    \draw[thick] (R3) -- node[lbl]{1} (R4);
    \draw[thick] (R4) -- node[lbl]{3} (R5);
    \draw[thick] (R5) -- node[lbl]{1} (R6);
    \draw[thick] (R6) -- node[lbl]{4} (R1);
    \draw[thick] (R2) -- node[lbl]{5} (R5);
    \draw[thick, dashed] (ATM1)  -- node[lbl]{1} (R1);
    \draw[thick, dashed] (BANK1) -- node[lbl]{1} (R4);
\end{tikzpicture}
\caption{Topologia de la red. Costos de enlace sobre cada arista. ATM1 y BANK1 son terminales (no routers).}
\end{figure}

La asignacion de nodos por participante es:

\begin{table}[H]
    \centering
    \caption{Asignacion de nodos por participante}
    \begin{tabular}{llll}
        \toprule
        \textbf{Participante} & \textbf{Nodos asignados} & \textbf{IP Tailscale} & \textbf{Puertos} \\
        \midrule
        Espana   & R1, R3      & 100.119.213.97 & 5001, 5003  \\
        Roberto  & R2, ATM1    & 100.84.155.75  & 5002, 6101  \\
        Tono     & R4          & 100.64.209.42  & 5004        \\
        Angel    & R5, R6, BANK1 & 100.71.208.101 & 5005, 5006, 6102 \\
        \bottomrule
    \end{tabular}
\end{table}

% ──────────────────────────────────────────
\section{Descripcion de los Algoritmos Utilizados y su Implementacion}

\subsection{Plano de Control: Protocolo Link-State}

\subsubsection{Descubrimiento de Vecinos (HELLO)}

Cada router envia periodicamente un paquete \texttt{HELLO} a todos sus vecinos
configurados inicialmente. Al recibir un \texttt{HELLO} de un nodo
desconocido, el router lo registra como vecino activo y actualiza su LSA. Si
un vecino deja de responder durante \texttt{dead\_interval} segundos, es
declarado caido y se propaga un LSA actualizado.

\begin{lstlisting}[language=Python, caption=Estructura del paquete HELLO]
{ "type": "HELLO", "from": "R1" }
\end{lstlisting}

\subsubsection{Publicacion del Estado del Enlace (LSA)}

Cuando el estado local cambia, el router construye un \textit{Link State
Advertisement} (LSA) con numero de secuencia incremental y lo inunda hacia
todos sus vecinos, salvo al remitente original (\textit{flooding controlado}).
Cada router ignora LSAs con numero de secuencia menor o igual al ultimo visto
del mismo origen.

\begin{lstlisting}[language=Python, caption=Estructura del paquete LSA]
{
    "type":   "LSA",
    "origin": "R1",
    "seq":    3,
    "links":  [{"to": "R2", "cost": 1}, {"to": "R3", "cost": 4}],
    "from":   "R1"
}
\end{lstlisting}

\subsubsection{Calculo de Rutas: Dijkstra}

Con la informacion de todos los LSAs recibidos, cada router construye el grafo
global de la red y aplica el \textbf{algoritmo de Dijkstra} para calcular la
ruta de costo minimo hacia cada destino. La implementacion usa una cola de
prioridad (\textit{min-heap}) para lograr complejidad $O((V + E) \log V)$.

\begin{table}[H]
    \centering
    \caption{Tabla de enrutamiento de R1 tras la convergencia}
    \begin{tabular}{llll}
        \toprule
        \textbf{destino} & \textbf{siguiente\_salto} & \textbf{costo} & \textbf{ruta} \\
        \midrule
        R2    & R2   & 1 & R1 $\to$ R2 \\
        R3    & R3   & 1 & R1 $\to$ R3 \\
        R4    & R2   & 2 & R1 $\to$ R2 $\to$ R3 $\to$ R4 \\
        R5    & R2   & 3 & R1 $\to$ R2 $\to$ R3 $\to$ R4 $\to$ R5 \\
        R6    & R6   & 4 & R1 $\to$ R6 \\
        BANK1 & R2   & 3 & R1 $\to$ R2 $\to$ R3 $\to$ R4 $\to$ BANK1 \\
        ATM1  & ATM1 & 1 & R1 $\to$ ATM1 \\
        \bottomrule
    \end{tabular}
\end{table}

La tabla se exporta automaticamente al archivo
\texttt{output/<nodo>\_nodo\_tabla\_enrutamiento.csv} cada vez que Dijkstra
recalcula las rutas, con columnas: \texttt{destino, siguiente\_salto, costo, ip, puerto}.

\subsection{Plano de Datos y Arquitectura de Capas}

\subsubsection{Reenvio salto a salto}

Al recibir un paquete de datos cuyo \texttt{nodo\_destino} no coincide con el
propio nodo, el router consulta su tabla de enrutamiento, obtiene la IP y
puerto del siguiente salto y reenvía el paquete por un nuevo socket TCP.
El formato de los paquetes de datos es:

\begin{lstlisting}[language=Python, caption=Paquete de datos]
{
    "nodo_origen":  "ATM1",
    "nodo_destino": "BANK1",
    "mensaje": {"action": "withdraw", "data": {"amount": 100}}
}
\end{lstlisting}

\subsubsection{Deteccion y Correccion de Errores: Hamming (7,4)}

La capa de datos aplica \textbf{Hamming (7,4)} a todos los bits del paquete
antes de transmitir. Los datos se dividen en bloques de 4 bits $[d_1, d_2,
d_3, d_4]$ y cada bloque se codifica en 7 bits $[p_1, p_2, d_1, p_3, d_2,
d_3, d_4]$ con tres bits de paridad calculados como:

\[
p_1 = d_1 \oplus d_2 \oplus d_4, \quad
p_2 = d_1 \oplus d_3 \oplus d_4, \quad
p_3 = d_2 \oplus d_3 \oplus d_4
\]

Esto permite detectar y corregir cualquier error de 1 bit por bloque. El
proceso de recepcion es:

\begin{enumerate}
    \item Recepcion de la cadena de bits codificada.
    \item Calculo del sindrome para detectar y corregir errores de 1 bit.
    \item Extraccion de los 4 bits de datos de cada bloque de 7.
    \item Deserializacion del mensaje JSON.
    \item Consulta de la tabla de enrutamiento.
    \item Reserializacion y recodificacion para el siguiente salto.
    \item Envio por socket.
\end{enumerate}

\subsubsection{Protocolo de Enmarcado Compartido}

Para garantizar la interoperabilidad entre grupos, se acordo el siguiente
formato de trama sobre TCP (terminada en salto de linea):

\[
\texttt{<algoritmo>|<cadena\_de\_bits>\textbackslash n}
\]

\begin{itemize}
    \item Plano de control (\texttt{HELLO}, \texttt{LSA}):
          \texttt{none|<bits>\textbackslash n} sin redundancia adicional.
    \item Plano de datos (paquetes ATM/BANK):
          \texttt{hamming|<bits\_codificados>\textbackslash n}.
\end{itemize}

La trama Hamming incluye un encabezado de 16 bits (\texttt{num\_blocks}) que
indica la cantidad de bloques de 7 bits, permitiendo al receptor delimitar
exactamente cuantos bits deben decodificarse:

\[
\underbrace{[b_{15}\ldots b_0]}_{\text{16-bit header}}
\;\Big|\;
\underbrace{[c_1^{(1)} \ldots c_7^{(1)}] \;\cdots\; [c_1^{(n)} \ldots c_7^{(n)}]}_{\text{n bloques Hamming de 7 bits}}
\]

\subsection{Cliente ATM y Servidor Bancario}

El nodo \texttt{ATM1} actua como \textbf{cliente}: el usuario interactua con
un menu de linea de comandos para ejecutar operaciones bancarias (\texttt{login},
\texttt{withdraw}, \texttt{logout}). El ATM envia los paquetes a su gateway
(\texttt{R1}) que los encamina a traves de la red hasta el banco.

El nodo \texttt{BANK1} actua como \textbf{servidor}: recibe y procesa las
solicitudes, administra sesiones y saldos, y envia respuestas de vuelta al ATM
a traves de la red de routers.

Ambos nodos no participan en el protocolo de enrutamiento: unicamente se
comunican con su gateway predeterminado mediante sockets.

% ──────────────────────────────────────────
\section{Resultados}

\subsection{Convergencia del Protocolo Link-State}

Tras iniciar todos los nodos, el protocolo converge en aproximadamente 2--3
segundos. Cada router imprime su tabla de enrutamiento al detectar cambios:

\begin{lstlisting}[caption=Salida de consola al iniciar el sistema]
==================================================
  PROPIETARIO: Espana  |  Modo: Tailscale
  Nodos: R1, R3
==================================================
Iniciando Router [R1] en 100.119.213.97:5001...
[R1] Publicando LSA seq=1 con 2 enlaces activos
[R1] HELLO -> R2 (ok)
[R1] Vecino detectado: R2
[R1] Tabla de ruteo actualizada (6 destinos):
   R2   -> next_hop=R2   cost=1 path=R1 -> R2
   R3   -> next_hop=R3   cost=1 path=R1 -> R3
   R4   -> next_hop=R2   cost=2 path=R1 -> R2 -> R3 -> R4
   BANK1 -> next_hop=R2  cost=3 path=R1 -> R2 -> R3 -> R4 -> BANK1
\end{lstlisting}

\subsection{Tablas de Enrutamiento Generadas}

Las tablas se exportan automaticamente al directorio \texttt{output/}. Ejemplo
del archivo \texttt{output/R1\_nodo\_tabla\_enrutamiento.csv}:

\begin{lstlisting}[caption=Contenido de R1\_nodo\_tabla\_enrutamiento.csv]
destino,siguiente_salto,costo,ip,puerto
R2,R2,1,100.84.155.75,5002
R3,R3,1,100.119.213.97,5003
R4,R2,2,100.84.155.75,5002
R5,R2,3,100.84.155.75,5002
R6,R6,4,100.119.213.97,5001
BANK1,R2,3,100.84.155.75,5002
ATM1,ATM1,1,100.84.155.75,6101
\end{lstlisting}

\subsection{Transaccion Bancaria End-to-End}

Se verifico la transaccion completa ATM1 $\to$ R1 $\to$ R2 $\to$ R3 $\to$
R4 $\to$ BANK1, con mensajes codificados en Hamming (7,4):

\begin{lstlisting}[caption=Flujo de reenvio observado en consola]
[R1][FWD] Reenviando a BANK1 via R2 (100.84.155.75:5002)
[R2][FWD] Reenviando a BANK1 via R3 (100.119.213.97:5003)
[R3][FWD] Reenviando a BANK1 via R4 (100.64.209.42:5004)
[R4][FWD] Reenviando a BANK1 via BANK1 (100.71.208.101:6102)
[BANK1] Peticion de ATM1: login
[BANK1] Respuesta enviada a ATM1: login_ok
\end{lstlisting}

\subsection{Pruebas Unitarias}

\begin{table}[H]
    \centering
    \caption{Resultados de la suite de pruebas}
    \begin{tabular}{lp{6.5cm}l}
        \toprule
        \textbf{Prueba} & \textbf{Descripcion} & \textbf{Resultado} \\
        \midrule
        \texttt{test\_hamming\_round\_trip}
            & Encode/decode Hamming (7,4) de cadena arbitraria. & PASS \\
        \texttt{test\_hamming\_frame\_with\_header}
            & Trama con cabecera de 16 bits. & PASS \\
        \texttt{test\_dijkstra\_shortest\_paths}
            & Rutas minimas sobre grafo de 4 nodos. & PASS \\
        \texttt{test\_framing\_encode\_decode}
            & Formato \texttt{hamming|bits\textbackslash n}. & PASS \\
        \texttt{test\_socket\_end\_to\_end}
            & Flujo login, withdraw y logout sobre sockets locales. & PASS \\
        \texttt{test\_interoperability}
            & Interoperabilidad con \texttt{Lab3-Redes}. & SKIP* \\
        \bottomrule
    \end{tabular}

    {\small * Se omite si el directorio \texttt{Lab3-Redes/} no esta presente.}
\end{table}

\noindent Resultado general: \texttt{Ran 6 tests in $\approx$2.7s --- OK (skipped=1)}.

% ──────────────────────────────────────────
\section{Discusion}

\paragraph{Convergencia y dinamismo.}
El protocolo Link-State con flooding de LSAs logra convergencia rapida: en la
topologia de 6 nodos, la estabilizacion ocurre en menos de 5 segundos tras el
inicio. La deteccion de vecinos caidos (basada en \texttt{dead\_interval}) y la
republica automatica del LSA permiten al sistema recuperarse ante fallos sin
intervencion manual.

\paragraph{Separacion de planos.}
La separacion entre plano de control y plano de datos fue clave para la
correcta implementacion. Los mensajes \texttt{HELLO} y \texttt{LSA} se
transmiten sin codificacion Hamming (algoritmo \texttt{none}) porque son
mensajes de senalizacion frecuentes donde la redundancia aumentaria el
overhead innecesariamente. El enunciado especifica que Hamming aplica
\textit{``una vez finalizadas las tablas de enrutamiento''}, lo que corresponde
exactamente al plano de datos.

\paragraph{Interoperabilidad.}
El acuerdo de un protocolo de trama compartido (\texttt{<algoritmo>|<bits>\textbackslash n})
fue el factor mas critico para la interoperabilidad entre grupos. Pequenas
diferencias en el formato (endianness de la cabecera de 16 bits, cantidad de
bits por bloque) habrian impedido la comunicacion. La sincronizacion de IPs
de Tailscale mediante \texttt{topology.json} y el generador de configuracion
del grupo companero resolvio la dependencia de configuracion estatica del otro
proyecto.

\paragraph{Limitaciones.}
El sistema actual no implementa autenticacion entre routers, por lo que un
nodo malicioso podria inyectar LSAs falsos. Tampoco se maneja la reconvergencia
ante particiones de red (el grafo se desconecta). Estas son limitaciones
esperables para el alcance del laboratorio.

% ──────────────────────────────────────────
\section{Conclusiones}

\begin{enumerate}[leftmargin=*]
    \item El protocolo Link-State permite a cada nodo construir una vision
          global de la red a partir de informacion distribuida localmente,
          calculando rutas optimas sin necesidad de un controlador centralizado.

    \item El uso de Tailscale simplifico drasticamente la conectividad entre
          maquinas en redes distintas, haciendo transparente la diferencia entre
          ejecutar nodos en una misma LAN o en internet.

    \item La separacion entre plano de control (sin Hamming) y plano de datos
          (con Hamming) optimiza el uso del ancho de banda: solo los datos de
          usuario cargan con la redundancia necesaria para la correccion de errores.

    \item La definicion de un protocolo de enmarcado compartido fue esencial
          para la interoperabilidad. Un estandar de interfaz bien definido
          permite integrar implementaciones independientes sin coordinacion
          adicional en tiempo de ejecucion.

    \item Las pruebas automatizadas fueron fundamentales para validar
          la correccion de cada componente de forma aislada antes de la
          integracion distribuida, reduciendo significativamente el tiempo
          de depuracion en red.
\end{enumerate}

% ──────────────────────────────────────────
\section*{Comentarios}

Durante el desarrollo se identifico que el mayor desafio no fue la
implementacion del algoritmo de Dijkstra en si, sino el manejo correcto de la
concurrencia: los hilos de HELLO, LSA periodico, deteccion de liveness y el
servidor TCP deben compartir el estado del router (grafo, tabla de enrutamiento,
secuencias vistas) con proteccion de mutex para evitar condiciones de carrera.

Tambien fue necesario garantizar que el flooding de LSAs no genere loops: cada
nodo debe registrar los numeros de secuencia vistos por origen y descartar LSAs
duplicados o desactualizados.

% ──────────────────────────────────────────
\begin{thebibliography}{9}

\bibitem{tanenbaum}
Tanenbaum, A. S., \& Wetherall, D. (2011).
\textit{Computer Networks} (5th ed.). Pearson.

\bibitem{kurose}
Kurose, J. F., \& Ross, K. W. (2021).
\textit{Computer Networking: A Top-Down Approach} (8th ed.). Pearson.

\bibitem{hamming}
Hamming, R. W. (1950).
Error detecting and error correcting codes.
\textit{Bell System Technical Journal}, 29(2), 147--160.

\bibitem{tailscale}
Tailscale Inc. (2024).
\textit{Tailscale Documentation}.
\url{https://tailscale.com/kb/}

\bibitem{dijkstra}
Dijkstra, E. W. (1959).
A note on two problems in connexion with graphs.
\textit{Numerische Mathematik}, 1(1), 269--271.

\end{thebibliography}

\end{document}
